\documentclass{article}
\usepackage{../math_template}

\author{Jonathan Kasongo}
\title{HMMT November General --- Q6/10}

\begin{document}
\maketitle

\begin{problem}{HMMT November General --- Q6/10}
Call a polygon \textit{normal} if it can be inscribed within a unit circle.
How many non-congruent normal polygons are there such that the square
of each side length is a positive integer?
\end{problem}

\begin{solution}{Solution}
The idea is that we can dissect any such polygon into a finite number of
triangles that look like so,

\begin{center}
\begin{asy}
/* Geogebra to Asymptote conversion, documentation at artofproblemsolving.com/Wiki go to User:Azjps/geogebra */
import graph; size(130);
real labelscalefactor = 0.5; /* changes label-to-point distance */
pen dps = linewidth(0.7) + fontsize(10); defaultpen(dps); /* default pen style */
pen dotstyle = black; /* point style */
real xmin = -0.148317663962672, xmax = 2.8495529768213816, ymin = -1.7315642208694166, ymax = 1.749754082680477;  /* image dimensions */
pen ududff = rgb(0.30196078431372547,0.30196078431372547,1.);
pair A = (0.,0.), B = (2.,-1.5), C = (2.,1.5);
pair m_AC = (1, 0.75), m_AB = (1, -0.75), m_BC = (2, 0);

filldraw(arc(A,0.26506675908800126,323.130102354156,396.869897645844)--(0.,0.)--cycle, ududff + opacity(0.30000001192092896) + ududff);
/* draw figures */
draw(A--C);
draw(C--B);
draw(B--A);
/* dots and labels */
//dot(A, dotstyle);
//dot(B,dotstyle);
//dot(C,dotstyle);
label("$1$", m_AC, NW * labelscalefactor);
label("$1$", m_AB, SW * labelscalefactor);
label("$\sqrt{a}$", m_BC, NE * labelscalefactor);

label("$\theta$", (0.3297286905226872,-0.04713626954620386), NE * labelscalefactor,ududff);
clip((xmin,ymin)--(xmin,ymax)--(xmax,ymax)--(xmax,ymin)--cycle);
/* end of picture */
\end{asy}
\end{center}

where $a$ is a positive integer. Now from Cosine law,
\[
\begin{aligned}
a &= 2(1 - \cos \theta) \\
\cos \theta &= 1 - \frac{a}{2} \\
\end{aligned}
\]
Since $\cos \theta \in [-1, 1]$ and $a \in \mathbb{N}$,
we can check that $a \in \{ 1, 2, 3, 4 \}$. Note that
$\theta \in [0, 180]$ since the angle is in a triangle. Now plugging in
each value of $a$ and solving yields
$$
\theta \in \{ 60, 90, 120, 180 \}.
$$
Since any normal polygon is just a collection of these triangles, we
are just looking for solutions to
\[
\begin{aligned}
60a + 90b + 120c + 180d &= 360\\
2a + 3b + 4c + 6d &= 12.
\end{aligned}
\]

\end{solution}

\end{document}
